\chapter{9}



Aus der Führerkabine klang Musik. Ein alter, amerikanischer Song lief.
Ein beleibter Mann beugte sich unter Anstrengung nach unten, als versuche
er etwas aufzuheben.

«Kann ich helfen, Signore?»

Der Mann schaute schräg zu Louis hoch. Über sein rotes Gesicht
flog augenblicklich ein sympathisches Lächeln. Die Luft war ein einziges
Gemisch aus Benzin und Öl.

«Gerne, junger Mann, mein Feuerzeug da, schauen Sie, unten, da, ich kann
es nicht erreichen, bin zu faul, um noch einmal auszusteigen.»

Nun lachte er laut und glucksend. Louis reichte es ihm.
Der Mann bedankte sich, stieg nun doch aus und streckte Louis seine Hand entgegen.

«Darf ich wissen, wohin Sie fahren, Signore?»

«Klar, Giovanni Colombo, mein Name. Und Sie, junger Mann?»

Ruckartig fühlte sich Louis überfahren. Himmel, er hatte sich seine neue
Geschichte erst rudimentär ausgedacht. Es musste sofort ein Name her.
Ihm war, als spreche ein anderer.

«Auch Giovanni, welch ein Zufall.»

Der Mann freute sich. Louis sah, dass ihm ein unterer Eckzahn fehlte.

\emph{«Meraviglioso»}, \footnote{«Wunderbar»} dann sind wir ja Namensvetter,» der Mann
klatschte in die Hände, «in die Schweiz fahre ich, ins Tessin, warum?»

Louis musste seine Überraschung verbergen. Er ahnte, dass ihm bereits
diese erste Begegnung die gewünschte Fahrgelegenheit bieten könnte.

«Ach, ich bin auch in die Schweiz unterwegs zu meiner Freundin und suche
eine Mitfahrgelegenheit. Natürlich kann ich auch was bezahlen.»

Er rief die Worte laut in den Motorenlärm der vorbeisausenden
Fahrzeuge.
Giovanni Colombo legte kurz den Kopf schräg, scannte Louis in raschem
Tempo von oben nach unten und zurück und reichte ihm erneut die Hand.
Wieder ertönte das herzliche Glucksen.

«Sicher, \emph{giovane Giovanni}, \footnote{«Junger Giovanni»} steig ein.»

Noch bevor sie losfuhren, wusste Louis, dass der Mann Grossvater war,
eine Enkelin hatte, die er immer mal wieder auf seine Fahrten mitnahm,
dass er sich freute, wenn er nicht alleine fahren musste, dass er gerne
alte, amerikanische Songs von Perry Como hörte, weil dessen Eltern auch
ursprünglich aus «Palena», einem kleinen Ort in den Abruzzen kamen wie
seine Ursprungsfamilie, dass Perry Comos Familie aber am Anfang des 20.
Jahrhunderts in die USA ausgewandert war und Como dort eine tolle
Karriere hingelegt hatte, und dass er, Giovanni Colombo, diese Musik
unendlich liebte, natürlich auch, weil sie ja eigentlich durch «Palena»
verbunden waren. Italiener halt, beide Familien.
Louis wurde vom vielen Nicken und Lächeln schier schwindlig. Er war
unendlich aufgekratzt und durfte
sich nichts anmerken lassen. Jetzt war er an der Reihe. Er hielt sich
kurz und setzte ein freundliches Gesicht auf.

«Wie schön, Giovanni.»

Es tat gut, italienisch zu sprechen.

«Ich bin Ihnen sehr dankbar, dass Sie mich mitnehmen, \emph{molto
grazie}.»

Er erzählte etwas von nachhaltigem Reisen und wiederholte die
Geschichte von der Freundin in der Schweiz.

«Es soll eine Überraschung werden,» Giovanni schien begeistert, «ich war
die halbe Nacht auf den Beinen und bin sehr müde. Wenn es Ihnen nichts
ausmacht, würde ich gerne ein bisschen schlafen.»

Louis blickte erwartungsvoll zu ihm hinüber.

«Aber natürlich, Gio-Gio, wenn ich dich so nennen darf.»

Er zwinkerte fröhlich mit dem linken Auge.

«Ich fahre mit Unterhaltung und hoffe, du hältst sie aus, schlaf gut.»

Giovanni tätschelte seinen Arm. Louis zog die Mütze weit über seine
Augen. Begleitet vom regelmässigen Motorengeräusch des Lasters, begann
er langsam dahin zu schlummern. Bis er von weit her Musik hörte.
Das war dieser Perry Como mit \emph{«Killing me softly»}. Louis kannte
den Song.  Es war lange her, seit er ihn bei seinem Grossvater gehört hatte. Diesmal verstand er auch die Worte.

\qrblock{Perry Como \emph{«Killing Me Softly»}}{media/QR-Codes/011_Killing_me_softly_Perry_Como_Spotify.png}

Von wegen, «softly» gekillt. Der abstruse Gedanke kam als erstes.
Giorgia hatte ihn knallhart abserviert.
Und doch war es zwecklos, sich gegen die Klänge und Comos Stimme zu
wehren. Ihm wurde warm. Die Situation erinnerte ihn an frühe
Kinderzeiten.  Er versuchte, die Tränen zurückzuhalten. Einige wenige rollten trotzdem
über seine Wangen. Giovanni schien es nicht zu bemerken.
Es war unfassbar. Giorgia, seine ihm über die Jahre ans Herz gewachsene
Giorgia, war einfach weg. Er fühlte eine Zuneigung, die ihn verwirrte. 
Er wollte nichts lieber als sie sanft halten und sie doch
gleichzeitig weit von sich stossen. Notwehr, das war’s, er hatte in
Notwehr gehandelt. Instinktiv hatte er ihren verbalen Angriff abgewehrt.
Der Gedanke gefiel ihm. Hätte er nicht postwendend Giorgias letzten
Blick vor sich auftauchen sehen, wäre er als Selbstrechtfertigung
durchgegangen.

\sectionbreak

Er konnte sich nicht erinnern, wann er vor seinem Zusammenbruch unter
der Brücke das letzte Mal geweint hatte. Es musste Jahre her sein. Neben
dem liebenswürdigen Mann fühlte er sich gerade geborgen und sicher.
Giovanni war ein bisschen wie Grandpère. Und er, Louis, der Kleine von
damals. Selbst wenn ihn der Gedanke beruhigte, kam er sich kindisch vor
und versuchte einmal mehr, Giorgias Blick zu verbannen.
Er schlief ein.

Der Nachrichten-Jingle weckte ihn Punkt acht. Louis lauschte den Worten
des Radiomoderators, ohne seine Stellung zu verändern. Giovanni sollte
weiterhin im Glauben sein, er schlafe. Seine Anspannung löste sich je
näher die Wetterprognose kam und verschwand bei den letzten Worten
« - und wünsche Ihnen einen sonnigen Tag, bis morgen, Ihr Claudio
Bruni».

Keine Fahndung also, weder nach ihm noch nach Giorgia.
Louis hob seine Brille und wischte sich über die Augen. Wie konnte er
nur annehmen, dass so rasch nach Giorgia oder sogar nach ihm gesucht
wurde. Wieder überlegte er sich mögliche Folgeszenen auf Giorgias Absturz.
Vielleicht war sie nur ein Stück gefallen, hängengeblieben und hatte
sich selbst retten können. Oder sie lag irgendwo und wurde noch gar nicht vermisst.
Oder sie war bereits zuhause und schlief ihren Rausch aus.
Die letzte Überlegung war die Unwahrscheinlichste. Sie diente einzig
seiner Hoffnung. Es war zum Verrücktwerden. Vielleicht sollte er sich mit Giovanni
unterhalten, um seine kreisenden Gedanken zu unterbrechen.

Mit dem konstanten Abbremsen des Lasters erübrigte sich eine
Entscheidung. Giovanni bog auf eine Raststätte ein und parkte das schwere
Gefährt zwischen zwei gigantischen Trucks. Louis spürte Giovannis Blick auf sich. Er begann sich
langsam zu strecken, schob die Kappe hoch und gähnte theatralisch. Mit
dem Öffnen seiner Augen schaute er direkt in Giovannis Gesicht. Der
strahlte.

«Guten Morgen Gio-Gio. Ich habe Hunger, du auch?»

«Oh, habe ich lange geschlafen? Wieviel Uhr ist denn?»

Louis setzte sich gerade auf.

«Bald halb neun, Zeit für ein gutes Frühstück. \emph{Andiamo?}\footnote{«Gehen wir?»}»

«Ja, eine gute Idee.»

Giovanni zog den Wagenschlüssel ab. Im Handschuhfach suchte er zwischen einem Durcheinander an 
Papieren nach seinem Geldbeutel und steckte ihn ein.

Sie stellten sich im Selbstbedienungsbereich Frühstücksteller zusammen.
Beim Anblick der verschiedenen Speisen schossen Louis Schuldgefühle
entgegen. Das «Ristorante». Marco. Er würde morgen nicht zur Arbeit
erscheinen. Marco würde sich fragen, wo Louis blieb und ohne Chef-Koch
dastehen. Das hatte sein langjähriger Arbeitgeber und Freund nicht
verdient. Wie konnte er sich nur bei ihm melden, ohne dass sein
Aufenthaltsort aufflog?

Unvermittelt fand sich Louis an der Kasse. Erleichtert, die Gedanken
vorläufig wegsperren zu müssen.
Er lud Giovanni ein. Der bedankte sich mit einer freundlichen Umarmung.
Sie setzten sich draussen an einen freien Tisch. Louis hatte ihn
ausgewählt. Er lag am Rande und leicht abseits der anderen Gäste.
Louis zog die Mütze tief ins Gesicht und sprach leise. Er würde dafür
sorgen, dass sie nicht lange hier sitzen blieben. Er stellte Giovanni
Fragen um Fragen.
Die Tische waren gut besetzt, die Morgentemperatur angenehm und der
Himmel über ihnen döste blau vor sich hin.
Giovanni war ein offenherziger Mensch. Louis war mehr als recht, dass er
ihn wenig Persönliches fragte. Giovanni schien zu gefallen, Louis sein
Leben zu erzählen. Mit seinen Geschichten kam der Mann ihm je länger je
näher. Er staunte, dass ihm Giovannis Familie irgendwie ans Herz wuchs.
Verrückt. Er kannte doch niemanden davon.
Der Mann war eine willkommene Ablenkung.
Schliesslich aber kam Louis nicht umhin, auf einige Fragen einzugehen.
Seine Antworten übertrafen sich an Bedeutungslosigkeit.

Der dichte Verkehr schien Giovanni zu fordern. Sie fuhren schweigend. Louis hing seinen Gedanken nach. 
Wider Erwarten hatten ihn Giovannis Erzählungen in vergangene Zeiten zurückgeschleudert. Die ganz persönlichen Dinge gingen keinen was an, wirklich? Erinnerungen an seinen Vater hatten ihm gerade noch gefehlt. 
Wie der nur dazu gekommen war, ihm diese «Verschwiegenheit» einzubläuen? Seine Mutter hatte auf ihre Art dagegengehalten. 
Louis staunte, wie plötzlich und präzise die damalige Szene vor ihm erschien. 

Es hatte sich über die Wochen eingespielt. 
Maman hatte an seinem Bett gesessen. Sie gehörte ganz ihm. Sie hielt seine Hand. 
Sie hörte zu, fragte zurück. Und manchmal, selten zwar, sprach sie von sich. Wie
sie die Dinge sah. Und jedes Mal, das hatte er deutlich gespürt, ging es
um etwas Bedeutungsvolles. Maman sah so vieles anders als Vater.
Eigentlich alles. Oft schloss sie mit ähnlichen Worten.

«Erwachsene Menschen haben ihre Meinungen. Papa ist dein Vater. Er sieht
die Dinge auf seine Art und ich auf meine. Das ist okay. Du weisst,
Louis, wir denken unterschiedlich. Darum ist es besser, dass wir nicht
mehr zusammenleben. Du bist ein kluger Junge, Louis. Du wirst deine
eigene Meinung finden.»

Bis an einem Sonntag. Als sie am Mittagstisch über diese
«Verschwiegenheit» gesprochen hatten, Maman und er. Schon alleine das
Wort schien sie zu reizen. Sie könne es nicht mehr hören. Sie sprach
überzeugter, heftiger. 
Louis erinnerte sich an seine Wut. Er hatte Vater verteidigt.

«Papa weiss das besser. Familiendinge darf man niemandem erzählen. Das
geht keinen was an.»

Grandpère hatte stumm daneben gesessen und war urplötzlich aufgestanden. Er
hatte sich am Kaffee verschluckt. Louis sah erschrocken zu ihm hoch. Ein
Gefühl von Traurigkeit und gleichzeitigem Trotz kam über ihn. Jetzt
waren sie seine Gegner. Zwei auf einmal. Louis war sich schwach
vorgekommen. Schon damals war es eigentlich klar. Grandpère und Vater mochten sich
nicht. Maman musste ihn und Lucia bewusst verschont haben. Vor all den
sich wiederholenden Tragödien, die sie mit ihrem Vater erlebte.

\sectionbreak

Louis schielte zu Giovanni hinüber. Auch er war ein Vater. Mit
Sicherheit ein anderer. 

«Wenn es rollt, sind wir schon bald an der Grenze, Gio-Gio.»

Louis atmete erstmal auf. Es war gut gelaufen. Seine Panik legte sich.
Niemand war auf sie aufmerksam geworden.
Und während Perry Comos \emph{«Magic Moments»} förmlich aus den
Musikboxen der Führerkabine sprudelten und sich Giovanni über die
Anwesenheit von Louis sichtlich freute, wähnte sich Louis für einen solch 
magischen Moment in einer fragilen Blase des Friedens.

\qrblock{Perry Como \emph{«Magic Moments»}}{media/QR-Codes/012_Magic_Moments_Perry_Como_Spotify.png}

Woher dieser kleine Lastwagenfahrer nur diese pure Fröhlichkeit nahm. Im
Takt hob und senkte Giovanni seine Schultern. Einmal links und einmal
rechts und sang einzelne Passagen treffsicher mit. Eingebettet in diese
Leichtigkeit und gepaart mit seiner lähmenden Erschöpfung stand Louis
kurz davor, Giovanni seinen Kummer anzuvertrauen.
Der Gedanke liess ihn zusammenzucken. So ein Blödsinn.
Giovannis weitere Fragen stoppte Louis mit Erklärungen rund um diese
nicht existierende Freundin und anderen Unwahrheiten. Er fühlte sich
schlecht dabei, ihn ein aufs andere Mal anzulügen.
Er musste unbedingt zu einer neuen Gitarre kommen. Erneut stellte er
sich schlafend.

Sie überquerten die Grenze in Chiasso. Giovanni steuerte die erste
Raststätte auf Schweizer Boden an.
